\documentclass[11pt,a4paper]{article}
\usepackage{fontspec,amsmath,amssymb,amsthm,microtype}
\setmainfont{lmroman10-regular.otf}[BoldFont=lmroman10-bold.otf,ItalicFont=lmroman10-italic.otf,BoldItalicFont=lmroman10-bolditalic.otf]
\usepackage{xcolor}
\usepackage[most]{tcolorbox}
\usepackage{mathrsfs}
\providecommand{\lt}{<}
\providecommand{\gt}{>}
\setmonofont{DejaVuSansMono.ttf}[BoldFont=DejaVuSansMono-Bold.ttf,ItalicFont=DejaVuSansMono-Oblique.ttf,BoldItalicFont=DejaVuSansMono-BoldOblique.ttf,Scale=0.86]
\newfontfamily\leanfallback{STIXTwoMath-Regular.otf}[Scale=0.86]
\definecolor{leanrule}{gray}{0.62}
\definecolor{leanback}{gray}{0.962}
\definecolor{leantitle}{gray}{0.22}
\definecolor{leanmuted}{gray}{0.38}
\newtcolorbox{leanbox}[1]{enhanced,breakable,colback=leanback,colframe=leanrule,
  boxrule=0.4pt,arc=0.6mm,left=2.4mm,right=2.4mm,top=1.8mm,bottom=1.8mm,
  before skip=0.8em,after skip=1em,fonttitle=\small\bfseries,coltitle=leantitle,
  colbacktitle=leanback,titlerule=0.3pt,title={#1}}
\newlength{\leancharwidth}
\newenvironment{leancode}{\par\smallskip\begingroup\ttfamily\footnotesize\raggedright
  \settowidth{\leancharwidth}{M}\setlength{\parskip}{0pt}\setlength{\parindent}{0pt}
  \frenchspacing}{\par\endgroup\smallskip}
\newcommand{\leanline}[2]{\hangindent=\dimexpr#1\leancharwidth+4\leancharwidth\relax
  \hangafter=1\noindent\hspace*{#1\leancharwidth}#2\par}
\newcommand{\leansource}[1]{\par{\footnotesize\color{leanmuted}#1}\par}
\newcommand{\leanlabel}[1]{\par\smallskip\noindent{\small\bfseries #1}\par}
\usepackage[margin=25mm]{geometry}
\usepackage[unicode,hidelinks]{hyperref}
\hypersetup{pdfauthor={LangevinLean Project}}
\setlength{\emergencystretch}{3em}
\setlength{\parskip}{0.4em}
\setlength{\parindent}{0pt}
\allowdisplaybreaks[2]
\raggedbottom
\title{\texorpdfstring{Retaining the actual initial coordinate through fresh updates}{Retaining the actual initial coordinate through fresh updates}}
\author{LangevinLean Project}
\date{}
\begin{document}\maketitle
Draft mathematical exposition; these notes may contain mistakes.


A sampled MALA update is identified with its explicit transition kernel. Extending the state by an unchanged initial coordinate yields an actual anchored recursion, whose law advances by the identity kernel on the anchor and the actual MALA kernel on the endpoint.


\section{\texorpdfstring{Actual sampled updates are kernel composition}{Actual sampled updates are kernel composition}}


\textbf{Claim.} Sampling an independent innovation and applying a measurable update is exactly composition with its sampled transition kernel. For the scalar MALA update, this kernel is the explicit Gaussian accepted/\allowbreak{}rejected kernel.



\textbf{Proof.} Product integration identifies the sampled Markov law with integration of the output kernel over the current state law. For MALA, integrate the independent uniform threshold exactly; its weighted Gaussian output agrees with the explicit transition formula.

\begin{leanbox}{Lean: \texttt{sampled\_\allowbreak{}markovStep\_\allowbreak{}eq\_\allowbreak{}comp}, \texttt{malaUpdate}, \texttt{malaUpdate\_\allowbreak{}measurable} and 1 more}
\leansource{\texttt{sampled\_\allowbreak{}markovStep\_\allowbreak{}eq\_\allowbreak{}comp} \quad WassersteinMixing/\allowbreak{}SemigroupRandomMALAAnchors.lean, lines 23–35}
\begin{leancode}
\leanline{0}{theorem\ sampled\_markovStep\_eq\_comp\ (F\ :\ A\ →\ Z\ →\ A)}
\leanline{4}{(hF\ :\ Measurable\ (Function.uncurry\ F))\ (η\ :\ Measure\ Z)\ [IsProbabilityMeasure\ η]}
\leanline{4}{(κ\ :\ Kernel\ A\ A)\ (hκ\ :\ ∀\ x,\ κ\ x\ =\ η.map\ (F\ x))}
\leanline{4}{(μ\ :\ Measure\ A)\ [IsProbabilityMeasure\ μ]\ :}
\leanline{4}{ULA.markovStep\ F\ η\ μ\ =\ κ\ ∘ₘ\ μ}
\end{leancode}
\leansource{\texttt{malaUpdate} \quad WassersteinMixing/\allowbreak{}SemigroupRandomMALAAnchors.lean, lines 43–44}
\begin{leancode}
\leanline{0}{def\ malaUpdate\ (V\ :\ E\ →\ ℝ)\ (h\ :\ ℝ)\ (x\ :\ E)\ (z\ :\ E\ ×\ I)\ :\ E\ :=}
\leanline{2}{(update\ V\ h\ (x,\ x)\ z).2}
\end{leancode}
\leansource{\texttt{malaUpdate\_\allowbreak{}measurable} \quad WassersteinMixing/\allowbreak{}SemigroupRandomMALAAnchors.lean, lines 46–50}
\begin{leancode}
\leanline{0}{theorem\ malaUpdate\_measurable\ (V\ :\ E\ →\ ℝ)\ (h\ :\ ℝ)\ (L\ :\ ℝ≥0)}
\leanline{4}{(hdiff\ :\ Differentiable\ ℝ\ V)\ (hLip\ :\ LipschitzWith\ L\ (gradient\ V))\ :}
\leanline{4}{Measurable\ (Function.uncurry\ (malaUpdate\ V\ h))}
\end{leancode}
\leansource{\texttt{malaUpdate\_\allowbreak{}law} \quad WassersteinMixing/\allowbreak{}SemigroupRandomMALAAnchors.lean, lines 53–69}
\begin{leancode}
\leanline{0}{theorem\ malaUpdate\_law\ (V\ :\ E\ →\ ℝ)\ (h\ :\ ℝ)\ (L\ :\ ℝ≥0)}
\leanline{4}{(hdiff\ :\ Differentiable\ ℝ\ V)\ (hLip\ :\ LipschitzWith\ L\ (gradient\ V))\ (x\ :\ E)\ :}
\leanline{4}{(innovationLaw\ (E\ :=\ E)).map\ (malaUpdate\ V\ h\ x)\ =}
\leanline{6}{MALAStationaryPair.explicitKernel\ V\ h\ x}
\end{leancode}
\leanlabel{How the Lean proof follows the argument}
\texttt{sampled\_\allowbreak{}markov\allowbreak{}Step\_\allowbreak{}eq\_\allowbreak{}comp} proves the generic equality by integration on measurable events. \texttt{mala\allowbreak{}Update} is the actual scalar uniform-threshold update; \texttt{mala\allowbreak{}Update\_\allowbreak{}law} applies the exact uniform-decision law to identify its kernel.
\end{leanbox}


Earlier results used here: \cite[Section~3.33.5]{LangevinLean2026ULASC}; \cite[Section~3.36.12]{LangevinLean2026ULASC}; \cite[Section~3.36.9]{LangevinLean2026ULASC}.

\section{\texorpdfstring{The actual anchored update and recursion}{The actual anchored update and recursion}}


\textbf{Claim.} Updating \((x_0,x)\) to \((x_0,\operatorname{MALA}(x,\xi,U))\) preserves the anchor and has kernel \(I\otimes K\). The resulting deterministic recursion equals the pair of its original initial value and the scalar sampled MALA recursion.



\textbf{Proof.} The first output coordinate is deterministic and the second has the proved explicit MALA law, so the sampled pair law is the Dirac anchor times that transition. Integrating over the current pair gives the parallel-kernel action. Inducting the deterministic recursion proves the exact anchor/\allowbreak{}endpoint identity.

\begin{leanbox}{Lean: \texttt{anchoredUpdate}, \texttt{anchoredUpdate\_\allowbreak{}measurable}, \texttt{anchoredUpdate\_\allowbreak{}law} and 2 more}
\leansource{\texttt{anchoredUpdate} \quad WassersteinMixing/\allowbreak{}SemigroupRandomMALAAnchors.lean, lines 72–73}
\begin{leancode}
\leanline{0}{def\ anchoredUpdate\ (V\ :\ E\ →\ ℝ)\ (h\ :\ ℝ)\ (p\ :\ E\ ×\ E)\ (z\ :\ E\ ×\ I)\ :\ E\ ×\ E\ :=}
\leanline{2}{(p.1,\ malaUpdate\ V\ h\ p.2\ z)}
\end{leancode}
\leansource{\texttt{anchoredUpdate\_\allowbreak{}measurable} \quad WassersteinMixing/\allowbreak{}SemigroupRandomMALAAnchors.lean, lines 75–79}
\begin{leancode}
\leanline{0}{theorem\ anchoredUpdate\_measurable\ (V\ :\ E\ →\ ℝ)\ (h\ :\ ℝ)\ (L\ :\ ℝ≥0)}
\leanline{4}{(hdiff\ :\ Differentiable\ ℝ\ V)\ (hLip\ :\ LipschitzWith\ L\ (gradient\ V))\ :}
\leanline{4}{Measurable\ (Function.uncurry\ (anchoredUpdate\ V\ h))}
\end{leancode}
\leansource{\texttt{anchoredUpdate\_\allowbreak{}law} \quad WassersteinMixing/\allowbreak{}SemigroupRandomMALAAnchors.lean, lines 81–91}
\begin{leancode}
\leanline{0}{theorem\ anchoredUpdate\_law\ (V\ :\ E\ →\ ℝ)\ (h\ :\ ℝ)\ (L\ :\ ℝ≥0)}
\leanline{4}{(hdiff\ :\ Differentiable\ ℝ\ V)\ (hLip\ :\ LipschitzWith\ L\ (gradient\ V))\ (p\ :\ E\ ×\ E)\ :}
\leanline{4}{(innovationLaw\ (E\ :=\ E)).map\ (anchoredUpdate\ V\ h\ p)\ =}
\leanline{6}{(Kernel.id\ {\leanfallback ∥}ₖ\ MALAStationaryPair.explicitKernel\ V\ h)\ p}
\end{leancode}
\leansource{\texttt{anchored\_\allowbreak{}markovStep} \quad WassersteinMixing/\allowbreak{}SemigroupRandomMALAAnchors.lean, lines 93–99}
\begin{leancode}
\leanline{0}{theorem\ anchored\_markovStep\ (V\ :\ E\ →\ ℝ)\ (h\ :\ ℝ)\ (L\ :\ ℝ≥0)}
\leanline{4}{(hdiff\ :\ Differentiable\ ℝ\ V)\ (hLip\ :\ LipschitzWith\ L\ (gradient\ V))}
\leanline{4}{(ρ\ :\ Measure\ (E\ ×\ E))\ [IsProbabilityMeasure\ ρ]\ :}
\leanline{4}{ULA.markovStep\ (anchoredUpdate\ V\ h)\ innovationLaw\ ρ\ =}
\leanline{6}{(Kernel.id\ {\leanfallback ∥}ₖ\ MALAStationaryPair.explicitKernel\ V\ h)\ ∘ₘ\ ρ}
\end{leancode}
\leansource{\texttt{anchored\_\allowbreak{}chain\_\allowbreak{}eq} \quad WassersteinMixing/\allowbreak{}SemigroupRandomMALAAnchors.lean, lines 103–112}
\begin{leancode}
\leanline{0}{theorem\ anchored\_chain\_eq\ (V\ :\ E\ →\ ℝ)\ (h\ :\ ℝ)\ (x\ :\ E)\ (w\ :\ ℕ\ →\ E\ ×\ I)\ (n\ :\ ℕ)\ :}
\leanline{4}{ULA.noiseDrivenChain\ (anchoredUpdate\ V\ h)\ (x,\ x)\ w\ n\ =}
\leanline{6}{(x,\ (ULA.noiseDrivenChain\ (update\ V\ h)\ (x,\ x)\ w\ n).2)}
\end{leancode}
\leanlabel{How the Lean proof follows the argument}
\texttt{anchored\allowbreak{}Update\_\allowbreak{}law} maps the exact scalar update law while keeping the first coordinate. \texttt{anchored\_\allowbreak{}markov\allowbreak{}Step} applies the generic sampled composition result. \texttt{anchored\_\allowbreak{}chain\_\allowbreak{}eq} is induction on the actual update recursion.
\end{leanbox}


Earlier results used here: \cite[Section~3.20.1]{LangevinLean2026ULASC}; \cite[Section~3.20.3]{LangevinLean2026ULASC}; \cite[Section~3.33.5]{LangevinLean2026ULASC}.

\begin{thebibliography}{ULA-LSI}
\bibitem[ULA-SC]{LangevinLean2026ULASC}
LangevinLean Project. \emph{Unadjusted Langevin under Strong Convexity}. LL-01, 2026. Draft report, content revision d34ead06987c. \url{https://odakyildiz.github.io/LangevinLean/docs/pdf/UnadjustedLangevin/FullProof.pdf}.
\end{thebibliography}

\end{document}
